\ifdefined\pdfinfoomitdate
  \pdfinfoomitdate=1
\fi
\ifdefined\pdftrailerid
  \pdftrailerid{}
\fi
\ifdefined\pdfsuppressptexinfo
  \pdfsuppressptexinfo=-1
\fi

\documentclass[11pt]{amsart}

\usepackage{amsmath,amssymb}
\usepackage{booktabs}
\usepackage{algorithm}
\usepackage{algpseudocode}
\usepackage[colorlinks=true,linkcolor=blue,urlcolor=blue]{hyperref}

\newcommand{\Q}{\mathbb Q}
\newcommand{\Z}{\mathbb Z}
\newcommand{\fB}{\mathcal B}
\newcommand{\rot}{\mathrm{rot}}
\newcommand{\code}[1]{\texttt{\detokenize{#1}}}
\newcommand{\algorithmcaption}[2]{%
  \refstepcounter{figure}\label{#2}%
  \begin{center}\small\textbf{Algorithm \thefigure.} #1\end{center}}

\title{Computational Details for Algorithmic Emerton--Gee Stacks}
\date{Version 14}

\begin{document}
\maketitle
\tableofcontents

\section{The bounded-cone and \texorpdfstring{$F_4$}{F4} computations}
\label{sec:f4}

Root data are generated from the $F_4$ Cartan matrix.  Roots are coefficient
tuples in Bourbaki order, and every rank calculation is exact.

\subsection{Field-degree bound}

For each of the $105$ downward-closed complements $K$, the program forms the
grouped cone matrix using every graph target.  A nonzero integral minor proves
its rank.  The dimension estimate is automatic for $d=[F:\Q_p]\ge2$.
Every prime divisor of the witness determinants belongs to $\{2,3\}$, so the
same ranks hold for $p>12$.

\begin{quote}
\begin{algorithmic}[1]
\Require The $F_4$ Cartan matrix.
\State Generate $\Phi^+$, its root order, and all $105$ complements $K$.
\For{each $K$}
  \State Form the grouped cone matrix with $\Psi_2=\Phi^+$.
  \State Find an integral maximal nonzero minor and factor its determinant.
  \State Solve the dimension inequality for $d$.
\EndFor
\State Return the degrees that still require a finite computation.
\end{algorithmic}
\end{quote}
\algorithmcaption{The field-degree bound}{alg:f4-degree-bound}

Only $d=1$ remains.

\subsection{Fine pairs without degree pruning}

At degree one, let $S\subset\Phi^+$ satisfy
\[
f(\alpha)=1\qquad(\alpha\in S)
\]
for some linear functional $f$.  Every such $S$ contains a linearly independent
subset with the same affine span.  The program starts with all independent
subsets and adjoins roots by breadth-first search.  Affine inconsistency is the
only reason that a branch stops.

The pair inventory has $S=\Psi_2\ne\emptyset$.  The empty subset satisfies the
affine equations, and both closure conditions hold for every
$\Psi_0\subseteq\Phi^+$.  For every nonempty complement
$C=\Phi^+\setminus\Phi_K$,
\[
\Psi_2=\emptyset,\qquad
\operatorname{obstruction}=0,\qquad
\varepsilon=\#C+\#(\Psi_0\setminus\Phi_K)\ge1.
\]
Indeed, the second summand is nonnegative and $\#C\ge1$.  Hence arbitrary
$\Psi_0$ gives no case with $\varepsilon=0$.  The program checks all $1048$
pairs $(C,r)$ with $C\ne\emptyset$ and
$0\le r=\#(\Psi_0\setminus\Phi_K)\le\#C$.

For each accepted $S$, choose $\alpha_0\in S$ and put
\[
L_S=\sum_{\alpha\in S}\Z(\alpha-\alpha_0),\qquad
\Psi_0=L_S\cap\Phi^+,
\qquad \Psi_2=S.
\]
If $\beta_1+\cdots+\beta_d\in\Psi_0$ with $d>0$ and
$\beta_i\in\Psi_2$, then
\[
d=f(\beta_1+\cdots+\beta_d)=0,
\]
which is impossible.  Thus the degree value is checked at every node but is
never used to stop the search.

\begin{quote}
\begin{algorithmic}[1]
\Require $\Phi^+$ and the reviewed pair list $P$.
\State For $S=\emptyset$, check
$\varepsilon=\#(\Phi^+\setminus\Phi_K)
+\#(\Psi_0\setminus\Phi_K)\ge1$.
\State Put every linearly independent nonempty subset of $\Phi^+$ in a queue.
\While{the queue is nonempty}
  \State Remove one subset $S$ and skip it if it was already examined.
  \State Check that $f(\alpha)=1$ for all $\alpha\in S$ is consistent.
  \State Compute $L_S$, $\Psi_0$, and both closure conditions.
  \State Check that no positive sum of roots in $S$ belongs to $\Psi_0$.
  \If{both closure conditions hold}
    \State Add $(\Psi_0,S)$ to the recomputed list $P'$.
  \EndIf
  \For{$\beta\in\Phi^+\setminus S$}
    \If{$S\cup\{\beta\}$ is affine-consistent}
      \State Add $S\cup\{\beta\}$ to the queue.
    \EndIf
  \EndFor
\EndWhile
\State Compare the ordered lists $P'$ and $P$ and their SHA-256 digests.
\end{algorithmic}
\end{quote}
\algorithmcaption{No-cutoff fine-pair computation}{alg:f4-pairs}

There are $9780$ independent seeds and $50347$ nonempty affine-consistent subsets.
The computation gives $4862$ fine pairs.  The recomputed and reviewed ordered
lists agree, with digest
\begin{center}
\scriptsize
\texttt{e1dc8ced6a02269877dbfbc4a908404657d8a415a3b503f54d429888af29d4be}.
\end{center}
All $3002$ relevant independent subsets occur among these pairs.

\subsection{Classification and the Borel count}

For nonempty $\Psi_2$, the cone calculation is applied to the $105$
possibilities for $K$ and the $4862$ fine pairs.  The empty complement accounts
for $4862$ entries.  Of the remaining $505648$ entries, $294388$ are excluded by
the number of graph-target rows and $211260$ receive polynomial Jacobian-rank
checks.  Exactly four have $\varepsilon=0$, namely Configurations I--IV.  The
separate empty-$\Psi_2$ calculation has none.

For their grouped cone matrices, the tuples
\[
(\#\text{variables},\#\text{relations},\operatorname{rank},\dim)
\]
are
\[
(28,17,17,11),\quad(29,17,17,12),\quad
(30,14,14,16),\quad(32,8,7,25).
\]
Edges with a common graph target contribute to one summed equation.

For split simple $F_4$, Configuration I has labels
\[
(a,1,0,1),\qquad a\in\Z/(p-1)\Z,
\]
and Configurations II--IV have labels
\[
(1,1,0,1),\qquad(0,1,0,1),\qquad(0,1,0,0).
\]
Consequently the number of configuration-labelled component records is
\[
(p-1)+3=p+2.
\]
At $p=17$, Configuration I contributes the $16$ records
$(\mathrm I,(a,1,0,1))$ for $0\le a\le15$, and the other three configurations
contribute one record each.  Hence the exact component count is
\[
16+3=19.
\]
There are $17$ distinct label tuples.  The tuple $(1,1,0,1)$ occurs for
$(\mathrm I,a=1)$ and $\mathrm{II}$, and $(0,1,0,1)$ occurs for
$(\mathrm I,a=0)$ and $\mathrm{III}$.

\section{Baker--Campbell--Hausdorff computation}
\label{sec:bch}

Let
\[
\operatorname{BCH}(X,Y)=\log(\exp(X)\exp(Y))=\sum_{m\ge1}Z_m,
\]
where $Z_m$ has bracket length $m$.  The implementation uses
\[
Z_1=X+Y
\]
and, for $m\ge2$,
\[
\begin{aligned}
mZ_m={}&\frac12[X-Y,Z_{m-1}]\\
&+\sum_{1\le q\le(m-1)/2}\frac{B_{2q}}{(2q)!}
\sum_{\substack{k_1+\cdots+k_{2q}=m-1\\k_i>0}}
[Z_{k_1},[\cdots,[Z_{k_{2q}},X+Y]]\cdots]].
\end{aligned}
\]
The first four terms are checked symbolically:
\[
\begin{aligned}
Z_1&=X+Y,\\
Z_2&=\tfrac12[X,Y],\\
Z_3&=\tfrac1{12}[X,[X,Y]]+\tfrac1{12}[[X,Y],Y],\\
Z_4&=\tfrac1{24}[X,[[X,Y],Y]].
\end{aligned}
\]

Independently of this recursion, put
\[
A(t)=\exp(tX)\exp(tY)-1
\]
in the free associative algebra and compute
\[
\log(1+A(t))\bmod t^{12}
=\sum_{r=1}^{11}\frac{(-1)^{r+1}}rA(t)^r\bmod t^{12}.
\]
For each $1\le m\le11$, the coefficient of $t^m$ is compared with the image of
$Z_m$ in the enveloping algebra of the free Lie algebra on $X,Y$.  All eleven
equalities hold.

Write
\[
F=\sum f_\alpha e_\alpha,\quad
G=\sum g_\alpha e_\alpha,\quad
P=\sum p_\alpha e_\alpha,\quad
Q=\sum q_\alpha e_\alpha.
\]
For $A_2$, with $[e_a,e_b]=e_c$, the $e_c$ coefficient of
\(\operatorname{BCH}(G,Q)-\operatorname{BCH}(F,P)\) is
\[
\frac12(g_aq_b-g_bq_a-f_ap_b+f_bp_a).
\]
For the $B_2$ test algebra
\[
[e_a,e_b]=e_c,\qquad[e_a,e_c]=2e_d,
\]
the $e_d$ coefficient is
\[
\begin{aligned}
&g_aq_c-g_cq_a
+\frac16(g_a-q_a)(g_aq_b-g_bq_a)\\
&\quad-f_ap_c+f_cp_a
-\frac16(f_a-p_a)(f_ap_b-f_bp_a).
\end{aligned}
\]
A numerical $B_2$ check gives obstruction coefficient $-4$ in $e_d$; adding
$\frac1{11}e_d$ to $F$ gives equality of the two BCH expressions.

The largest positive-root height in $F_4$ is $11$.  The program constructs the
positive Chevalley Lie algebra and computes all eleven homogeneous terms.
It independently compares every F4 structure constant and every leading
root-height Jacobian with \code{roots.py} and \code{algebra.py}.  For Configurations
I--IV, a separate root-sum check proves that a term with one $\Psi_0$
perturbation and at least two $\Psi_2$ factors cannot contribute to a
$\Psi_2$ first-order rotation equation.

\section{Bipartite root-poset computations}
\label{sec:root-posets}

For a bipartite graph, the reduction procedure repeatedly removes a graph
source of out-degree one together with its unique adjacent graph target;
graph sources of out-degree zero are discarded.  If a graph remains, its
Jacobian has entries
\[
J_{\alpha,\beta}=N_{\beta,\alpha-\beta}x_{\alpha-\beta}.
\]
For $r$ graph-target rows and $c\ge r$ graph-source columns, the program tests
the $\binom cr$ column subsets until it finds a nonzero determinant consisting
of one monomial.

\subsection{Exceptional Chevalley types}

The finite graph computation is as follows.

\begin{quote}
\begin{algorithmic}[1]
\Require A root system $\Phi$ of exceptional Chevalley type.
\State Construct the directed root poset.
\For{each height $h$}
  \For{each graph-target subset $T$ at height $h+1$}
    \State Keep the height-$h$ graph sources and the edges ending in $T$.
    \State Repeatedly remove degree-one graph sources and adjacent targets.
    \If{no edge remains}
      \State Record success and continue.
    \EndIf
    \State Form $J_{\alpha,\beta}$ and test graph-source column subsets.
    \State Fail unless one maximal determinant is a nonzero monomial.
  \EndFor
\EndFor
\end{algorithmic}
\end{quote}
\algorithmcaption{Absolute root-poset check}{alg:absolute-root-posets}

The numerical summary is
\[
\begin{array}{c|rr}
\toprule
\text{Type}&\text{graph-target selections}&\text{failures}\\
\midrule
F_4&48&0\\
E_6&138&0\\
E_7&384&0\\
E_8&1196&0\\
\bottomrule
\end{array}
\]
and the nontrivial determinants, up to signs from row and column order, are
\begingroup
\small
\[
\begin{array}{c|c|l}
\toprule
\text{Type}&h&\det J\\
\midrule
F_4&3&3x_{0001}x_{0010}x_{1000}\\
E_6&2&2x_{000010}x_{001000}x_{010000}\\
E_6&3&3x_{000001}x_{000010}x_{001000}x_{010000}x_{100000}\\
E_7&2&-2x_{0000100}x_{0010000}x_{0100000}\\
E_7&3&-3x_{0000010}x_{0000100}x_{0010000}x_{0100000}x_{1000000}\\
E_7&4&2x_{0000001}x_{0001000}x_{0010000}x_{0100000}\\
E_7&8&-2x_{0000001}x_{0000010}x_{0000100}x_{0010000}\\
E_8&2&-2x_{00001000}x_{00100000}x_{01000000}\\
E_8&3&-3x_{00000100}x_{00001000}x_{00100000}x_{01000000}x_{10000000}\\
E_8&4&2x_{00000010}x_{00010000}x_{00100000}x_{01000000}\\
E_8&5&-5x_{00000001}x_{00000010}x_{00000100}x_{00010000}x_{00100000}x_{01000000}x_{10000000}\\
E_8&8&-2x_{00000010}x_{00000100}x_{00001000}x_{00100000}\\
E_8&9&3x_{00000001}x_{00000100}x_{00001000}x_{00010000}x_{00100000}x_{10000000}\\
E_8&14&2x_{00000010}x_{00001000}x_{00010000}x_{01000000}\\
\bottomrule
\end{array}
\]
\endgroup

\subsection{Type \texorpdfstring{$D_n$}{Dn}}

Use the Bourbaki simple roots
\[
\alpha_i=e_i-e_{i+1}\ (1\le i\le n-1),\qquad
\alpha_n=e_{n-1}+e_n.
\]
The implementation converts the positive roots by
\[
e_i-e_j=\sum_{k=i}^{j-1}\alpha_k
\]
and
\[
e_i+e_j=
\begin{cases}
\displaystyle\sum_{k=i}^{n-2}\alpha_k+\alpha_n,&j=n,\\[4pt]
\displaystyle\sum_{k=i}^{j-1}\alpha_k
+2\sum_{k=j}^{n-2}\alpha_k+\alpha_{n-1}+\alpha_n,&j<n.
\end{cases}
\]

For $h=2m$ with $1\le m\le\lfloor(n-2)/2\rfloor$, put $a=n-2m$ and define
\[
\begin{aligned}
s_r&=e_{a+r}+e_{n-r} &&(0\le r<m),\\
t_r&=e_{a-1+r}+e_{n-r} &&(0\le r\le m),\\
u&=e_a-e_n,\qquad
v=e_{a-1}-e_{n-1},\qquad
w=e_{a-1}-e_n.
\end{aligned}
\]
After the graph reduction above, the remaining edges are exactly
\[
s_r\longrightarrow t_r,t_{r+1},\qquad
u\longrightarrow w,t_1,\qquad
v\longrightarrow w,t_0.
\]
This graph is $\vec C_6$ joined at one graph target to the alternating path
$\vec P_{2m-1}$.  At every other height no edge remains.  For $D_8$ at height
$6$, one has $m=3$ and $\vec C_6\cup_{K_1}\vec P_5$.

For every subset of graph targets, a nonempty graph remains precisely when the
subset contains all graph targets in the graph above; in that case the
remaining graph is the one just written.  Its square Jacobian satisfies
\[
\det J_{n,m}=\pm2x_{\alpha_{n-2m-1}}
\prod_{i=n-m}^{n}x_{\alpha_i}.
\]
For $D_8$ at height $6$, the graph-source and graph-target strings are
\[
\begin{array}{c|c}
00011211&00012211,\ 00111211\\
00111111&00111211,\ 01111111\\
01111101&01111111,\ 11111101\\
01111110&01111111,\ 11111110\\
11111100&11111101,\ 11111110
\end{array}
\]
and
\[
\det J_{8,3}=2x_{00000001}x_{00000010}x_{00000100}
x_{00001000}x_{10000000}.
\]

The finite verification has the following ranges and counts:
\[
\begin{array}{c|c|r}
\toprule
\text{check}&\text{ranks}&\text{number examined}\\
\midrule
\text{all adjacent-height graphs}&4\le n\le32&928\\
\text{all graph-target subsets}&4\le n\le12&15888\\
\text{determinants}&4\le n\le12&25\\
\bottomrule
\end{array}
\]
There are $225$ nonempty graphs in the first row, $677$ nonempty selections in
the second row, and the determinant coefficients are exactly $-2$ and $2$.

\subsection{Non-Borel root posets}

The non-Borel test first applies the following disjunction: a relative cycle is
non-obstructing, or its associated absolute graph passes
Algorithm~\ref{alg:absolute-root-posets}.  The non-obstructing condition is
therefore checked before an absolute Jacobian is formed.

\begin{quote}
\begin{algorithmic}[1]
\Require A Levi subgroup, an elliptic Weyl group element $w$, and its relative root orbits.
\State Construct the $w$-relative root poset.
\State Remove individually non-obstructing graph targets.
\State Mark graph-target subsets with no non-obstructing pair.
\For{each adjacent-height bipartite graph}
  \State Remove degree-one graph sources and adjacent targets repeatedly.
  \For{each marked subset}
    \State Restrict to that subset and repeat the graph reduction.
    \State Replace each relative vertex by all absolute roots in its $w$-orbit.
    \State Add an edge when target minus source is a positive root.
    \State Apply Algorithm~\ref{alg:absolute-root-posets} if an edge remains.
  \EndFor
\EndFor
\end{algorithmic}
\end{quote}
\algorithmcaption{Non-Borel cycle check}{alg:non-borel}

The computation generates all Levi and elliptic Weyl group cases for the
exceptional Chevalley types and gives
\[
\begin{array}{c|rr}
\toprule
\text{Type}&\text{cases}&\text{failures}\\
\midrule
F_4&19&0\\
E_6&68&0\\
E_7&145&0\\
E_8&304&0\\
\bottomrule
\end{array}
\]

\section{Characteristic-uniform rotation checks}
\label{sec:rotation}

The rotation computation is performed over $\Q$.  It records integral data
whose reductions give a smooth point in every characteristic $p>12$.
For a rotation root $\delta_{\rot}$ and $\alpha\in\Psi_2$, the program follows
the root string $\alpha+n\delta_{\rot}$.  At step $n$ it multiplies the current
coefficient by
\[
N_{-\delta_{\rot},\alpha+n\delta_{\rot}}/n.
\]
For Configuration I this gives
\[
c_{0100}\longmapsto c_{0100}+2\mu c_{0110}-\mu^2c_{0120},\qquad
c_{0001}\longmapsto c_{0001}-\mu c_{0011}.
\]

\begin{quote}
\begin{algorithmic}[1]
\Require $(K,\Psi_0,\Psi_2)$ and a rotation root $\delta_{\rot}$.
\State Generate the Chevalley constants and root-string substitutions.
\State Form the divided first-order deformation equations over $\Q$.
\State Check from root sums that higher BCH words do not enter these equations.
\For{the deterministic sequence of specializations and $\Q$-points}
  \State Evaluate every equation and the square Jacobian determinant exactly.
  \If{all equations vanish and $\det J\ne0$}
    \State Factor the denominators and $\det J$.
    \State Check every required nonzero coordinate.
    \State Stop when every prime divisor is at most $12$.
  \EndIf
\EndFor
\State Fail if no such $\Q$-point was found.
\end{algorithmic}
\end{quote}
\algorithmcaption{Rotation check over $\Q$}{alg:rotation}

The exact output is
\[
\begin{array}{c|c|r|c}
\toprule
\text{Configuration}&\delta_{\rot}&\det J&\text{prime divisors}\\
\midrule
\mathrm{I}&0010&432&2,3\\
\mathrm{II}&0010&704&2,3,5,11\\
\mathrm{III}&1000&15360&2,3,5\\
\mathrm{IV}&0010&895795200&2,3,5\\
\bottomrule
\end{array}
\]
Every equation vanishes at its recorded $\Q$-point.  After inverting the
listed primes, that point defines a smooth section; reduction gives a smooth
point for every $p>12$.

The program also checks the dependence on the rotation root without changing
the four choices above.  It evaluates precisely
\[
\begin{array}{c|c|r|c|c}
\toprule
\text{Configuration}&\delta_{\rot}&\det J&\text{prime divisors}&\text{result}\\
\midrule
\mathrm{IV}&0010&895795200&2,3,5&\text{success}\\
\mathrm{IV}&0011&-30800&2,3,5,7,11&\text{success}\\
\mathrm{IV}&1000&315/4&2,3,5,7&\text{success}\\
\mathrm{III}&1000&15360&2,3,5&\text{success}\\
\mathrm{III}&0010&-29160&2,3,5&\text{success}\\
\bottomrule
\end{array}
\]
For every row, the recorded $\Q$-values give
\[
f_1=\cdots=f_n=0,\qquad \det J\ne0,\qquad
\prod_{\alpha\in\Psi_2}a_{t,\alpha}a_{p,\alpha}\ne0,
\]
and every prime divisor is at most $12$.  Hence all five roots work for every
$p>12$; among these choices the successful root is not unique for either
Configuration III or Configuration IV.  In particular, none of the five has
no solution, an everywhere-zero Jacobian determinant on its solution set, or
a failure of the required nonvanishing condition.

\section{Implementation and data}
\label{sec:implementation}

The implementation and reviewed data are available at
\begin{center}
\url{https://github.com/mocham/AlgEG/tree/grob-v14/code/v4}.
\end{center}
The commands, runtimes, deterministic JSON format, checksums, and module list
are recorded in
\begin{center}
\url{https://github.com/mocham/AlgEG/blob/grob-v14/code/README.md}.
\end{center}

The CLI writes to \code{code/v4/DATA} unless \code{--output-directory} gives
another directory.  With \code{--check}, it disables Python bytecode before local imports,
then recomputes and compares every requested JSON and SHA-256 file without
writing.  The full non-writing run is
\[
\code{python cli.py all --check --output-directory DATA}.
\]

\end{document}
